\section{FLORES-200: primero, la infraestructura de evaluación}

Una vez establecida la necesidad real, el siguiente paso no es entrenar de inmediato, sino responder primero «cómo saber si el sistema de verdad mejora». Un problema frecuente en la investigación con pocos recursos es que faltan el conjunto de entrenamiento, el conjunto de validación y las métricas; sin un conjunto de evaluación común, comparable y traducido por personas, las mejoras del modelo no pueden verificarse entre idiomas. La serie FLORES se creó para cubrir esta carencia de infraestructura.

\subsection{El orden de investigación Evaluation-first}

Esta sección explica por qué la clase aborda primero la evaluation y después el training. La clase coloca de forma explícita los evaluation datasets antes que los training datasets. Este orden expresa una disciplina de ingeniería: si aún no se han definido las condiciones de éxito, la distribución de prueba y los tipos de falla, las iteraciones de entrenamiento solo optimizarán el objetivo sustituto que resulte más cómodo de medir.

\lecturefigure{slide-09-evaluation-datasets.jpg}{NLLB trata primero los datos de evaluación y después los datos de entrenamiento}{Video de la clase de Stanford 00:07:24}

\begin{knowledgebox}{Lectura de la figura: la portada de sección también expresa prioridades}
Esta slide casi no tiene detalles técnicos, pero muestra la lógica con la que se organiza el curso: el benchmark no es un accesorio que se agrega cuando el modelo ya está terminado, sino el sistema de coordenadas que orienta el entrenamiento. FLORES-200 define primero la comparabilidad entre idiomas; solo así el data mining, el MoE y el curriculum posteriores pueden juzgar aciertos y errores en la misma escala.
\end{knowledgebox}

\teachervoice{Fan dice que quiere «hablar primero de los evaluation data, porque hacer bien la evaluation es sumamente importante». Esta frase debe leerse como un método de investigación: cuando no se tiene un conjunto de prueba confiable, no hay que tomar la loss de entrenamiento ni los resultados de unos cuantos idiomas con muchos recursos como representativos del conjunto.}

\subsection{De FLORES a FLORES-200}

A continuación hay que entender cómo el benchmark se amplió de una generación a otra. FLORES se construyó en un principio en torno a la evaluación de unos pocos idiomas con pocos recursos; después se amplió a FLORES-101, y NLLB llevó la cobertura hasta FLORES-200. En clase se mencionó que el nombre proviene de Facebook Low Resource y que, aunque la empresa cambió su nombre a Meta, el proyecto conservó el nombre FLORES.

\lecturefigure{slide-10-flores-101.jpg}{FLORES-101 sentó el precedente para la evaluación a gran escala con pocos recursos}{Video de la clase de Stanford 00:08:12}

\begin{knowledgebox}{Lectura de la figura: el valor del benchmark proviene de oraciones fuente comunes y de la traducción humana}
FLORES hace que muchos idiomas compartan un mismo grupo de oraciones con correspondencia semántica, lo que permite comparar la traducción into-English, out-of-English y non-English. No consiste en recopilar páginas web existentes y dividirlas al azar, sino en producir puntos de referencia de evaluación más estables mediante traducción profesional y procesos de control de calidad.
\end{knowledgebox}

\begin{importantbox}{Por qué las direcciones de traducción crecen de forma combinatoria}
Si hay $N$ idiomas y se consideran $A\rightarrow B$ y $B\rightarrow A$ como direcciones distintas, existen como máximo
\[
D=N(N-1)
\]
direcciones de traducción dirigidas. $N$ representa el número de idiomas y $D$, el número de direcciones sin contar la copia al mismo idioma. Para $N=200$, $D=39{,}800$; esto explica las más de cuarenta mil direcciones que se mencionaron en clase, y también por qué es imposible que los datos de entrenamiento sean igual de abundantes para cada dirección.
\end{importantbox}

\subsection{Los cuatro atributos de diseño de FLORES}

Esta sección desglosa «tener un conjunto de evaluación» en decisiones de diseño concretas. La clase muestra las decisiones de diseño de FLORES con progressive bullets, cuyo estado final incluye cuatro puntos: centrarse en idiomas con pocos recursos, admitir many-to-many, abarcar temas y dominios diversos, y conservar el document-level context. Cada uno de estos atributos responde, respectivamente, al sesgo de cobertura, al sesgo de dirección, al sesgo de dominio y al sesgo de oraciones aisladas.

\lecturefigure{slide-11-flores-properties.jpg}{El diseño de FLORES-200: pocos recursos, muchos a muchos, varios dominios y nivel de documento}{Video de la clase de Stanford 00:09:57}

\begin{knowledgebox}{Lectura de la figura: cuatro atributos frente a cuatro distorsiones del benchmark}
Elegir solo idiomas con muchos recursos sobrestima la madurez del sistema; evaluar solo English-centric directions oculta la degradación en las direcciones non-English; evaluar un solo dominio hace pasar la coincidencia de dominio por una capacidad general; evaluar solo oraciones aisladas pasa por alto los pronombres, la coherencia y las ambigüedades que dependen del contexto. FLORES no puede eliminar todos los sesgos, pero reduce notablemente estas distorsiones y permite discutirlas.
\end{knowledgebox}

\begin{warningbox}{Many-to-many evaluation no equivale a many-to-many supervision}
Que el conjunto de evaluación permita construir pares de comparación para todos los idiomas no significa que en la etapa de entrenamiento se disponga de corpus paralelos para todas las direcciones. En la sesión de Q\&A de la clase se dijo explícitamente que NLLB no extrajo las $200\times200$ direcciones completas; muchas direcciones dependen de la zero-shot transfer que surge de una representación multilingüe compartida.
\end{warningbox}

\subsection{Traducción profesional, revisión automática y revisión humana}

El proceso de producción de FLORES comienza con el acuerdo sobre los estándares lingüísticos; después, los documentos se entregan a los traductores y pasan por revisiones automáticas, un reviewer humano y post-editing. Si la calidad queda por debajo del umbral, la muestra regresa a la etapa de corrección. Este ciclo muestra que los datos de evaluación de alta calidad no se obtienen con una sola compra, sino que requieren la acción conjunta de estándares lingüísticos, herramientas de revisión y juicio profesional.

\lecturefigure{slide-12-flores-collection-pipeline.jpg}{El proceso de FLORES, desde el acuerdo sobre estándares lingüísticos hasta la revisión y la posedición}{Video de la clase de Stanford 00:11:30}

\begin{knowledgebox}{Lectura de la figura: los ciclos entre flechas importan más que el nodo final de término}
Primero se fijan los estándares de script, ortografía y variante, para que los traductores sepan qué forma debe adoptar el texto de destino; las revisiones automáticas detectan omisiones de traducción, problemas de formato y caracteres anómalos, pero no pueden juzgar por sí solas la naturalidad semántica; el reviewer y el post-editing se encargan de regresar las muestras que no pasan. El `Quality > 90%` de la figura es la condición que se mostró en clase para este proceso y no debe extrapolarse como un umbral fijo de control para todos los proyectos lingüísticos.
\end{knowledgebox}

\teachervoice{El docente enfatiza que el primer paso, y el verdaderamente difícil, suele ser la alignment on language standards. Si una comunidad usa al mismo tiempo varios sistemas de escritura, ortografías y variantes locales, el equipo debe decidir primero qué se evalúa; no puede esperar a que termine el entrenamiento del modelo para completar las etiquetas.}

\subsection{Traductores, estándares, sistemas de escritura y variantes locales}

Por último, volvemos a las restricciones de producción del conjunto de evaluación. La última slide sobre datos de evaluación resume las dificultades de recopilación en tres tipos: encontrar traductores adecuados, coordinar los estándares lingüísticos y manejar la diversidad de sistemas de escritura y las variantes locales. Parecen asuntos de «operación de datos», pero en realidad determinan directamente si el benchmark mide el idioma de destino o solo una forma de escritura institucionalizada.

\lecturefigure{slide-13-flores-collection-challenges.jpg}{Los retos de traductores y estándares lingüísticos en la recopilación de datos de FLORES}{Video de la clase de Stanford 00:12:54}

\begin{knowledgebox}{Lectura de la figura: cada bullet se propaga al modelo y a las métricas}
Un reclutamiento insuficiente de traductores amplía la varianza; una mala elección de estándar hace que variantes legítimas se califiquen como errores; la mezcla de sistemas de escritura afecta la tokenización y la language identification; la ausencia de variantes locales hace que el modelo obtenga buenos resultados en el benchmark sin poder atender a las comunidades reales. El estándar de datos no es un recipiente neutral, sino el measurement model de todo el sistema.
\end{knowledgebox}

\begin{warningbox}{Translationese y English-centric source}
Un conjunto de evaluación traducido por personas aún puede contener translationese, es decir, oraciones de destino que conservan la estructura del idioma fuente; si el contenido fuente proviene sobre todo de un contexto en inglés, también puede carecer de temas culturales locales. La traducción humana de alta calidad mejora notablemente la comparabilidad, pero no equivale de forma automática a la distribución de un corpus nativo.
\end{warningbox}

\subsection{Resumen de la sección}

FLORES-200 establece primero un sistema de coordenadas de evaluación común para doscientos idiomas y, sobre esa base, hace comparable la investigación de modelos. Su valor proviene de la cobertura de idiomas con pocos recursos, el diseño muchos a muchos, la variedad de dominios, el contexto de documento y el proceso de traducción profesional; sus limitaciones incluyen la elección de estándares, el translationese y el sesgo del contenido fuente.
